% problem-000179 1 主族元素反应方程式
\section*{1. 主族元素反应方程式}
⽔源洁净事关⼈⺠的⽣命健康。重⾦属离⼦、微⽣物含量等多种指标是⽔质考察的重点，锰含量就是其中⼀项。我国⽣活饮⽤⽔卫⽣标准规定锰含量不得超过0.1 mg L−1。锰含量常⽤分光光度法检测，实验过程如下：向⽔样中加⼊甲醛肟溶液，再加⼊适量氢氧化钠，放置20分钟，若⽔样中含锰（通常为 $\mathrm { M n ^ { 2 + } }$ ，其配合物⼏乎⽆⾊），会产⽣棕⾊配合物（锰离⼦与甲醛肟配⽐为1:6），测量其吸光度。

\subsection*{1-1}
甲醛肟(CH_{2}NOH)可以由盐酸羟胺与甲醛按1:1在⽔溶液中制得，写出反应⽅程式。

\subsection*{1-2}
写出棕⾊配合物形成的化学反应⽅程式。

\subsection*{1-3}
利⽤此⽅法测定浓度为 $8 . 0 0 { \times } 1 0 ^ { - 5 } \mathrm { m o l { \cdot } L ^ { - 1 } }$ 的标液，⽐⾊⽫宽1 cm，吸光度为0.880。从⽔源地取样，

将1 L⽔样浓缩⾄10 mL，若采⽤⽰差分光光度法，以 $1 . 6 0 \times 1 0 ^ { - 4 } \mathrm { m o l } \mathrm { L } ^ { - 1 }$ 锰标液作参⽐，测得样品的吸光度为0.200，计算⽔样中锰的含量，判断此⽔样是否合格。
